\documentclass[a4paper,11pt]{article}
\usepackage[utf8]{inputenc}
\usepackage[T1]{fontenc}
\usepackage[french]{babel}
\usepackage{lmodern}
\usepackage{amsmath,amssymb,amsthm}
\usepackage{mathrsfs}
\usepackage{enumitem}
\usepackage{multicol}

\setlist[enumerate]{label=\textbf{\textcolor{Couleur8}{\arabic*.}}, left=1em}

% Si vous voulez aussi modifier les listes enumerate imbriquées :
\setlist[enumerate,2]{label=\textcolor{Couleur8}{\alph*)}}
\setlist[enumerate,3]{label=\textcolor{Couleur8}{\roman*)}}
\setlist[enumerate,4]{label=\textcolor{Couleur8}{\Alph*)}}
\usepackage{graphicx}
\usepackage{xcolor}
\usepackage{tcolorbox}
\usepackage{geometry}
\usepackage{fancyhdr}
\usepackage{titlesec}
\usepackage{cancel}
\usepackage{ifthen}
\usepackage{tikz}
\usetikzlibrary{arrows.meta}
\usepackage{tkz-tab}
\usepackage{eurosym}

% OPTION PROF/ÉLÈVE
% Décommentez UNE des deux lignes suivantes :
\newboolean{prof}\setboolean{prof}{true}   % Version professeur (avec corrections des devoirs)
\newboolean{prof}\setboolean{prof}{true}  % Version élève (sans corrections des devoirs)

\definecolor{Couleur1}{HTML}{4A5D7A} %Th
\definecolor{Couleur2}{HTML}{A5B5D6} %Prop
\definecolor{Couleur3}{HTML}{A8C68A} %Méthode
\definecolor{Couleur6}{HTML}{8A9CC1} %Def
\definecolor{Couleur7}{HTML}{E6B459} %Rq Voc
\definecolor{Couleur8}{HTML}{D36740} %Titres
\definecolor{correction}{HTML}{D36740} % Couleur pour les corrections
\definecolor{coopmathscorr}{HTML}{F15929} % Couleur corrections coopmaths

% Configuration des marges
\geometry{
	top=2cm,
	bottom=2cm,
	inner=1.5cm,
	outer=1.5cm,
}

% Police
\renewcommand{\familydefault}{\sfdefault}

% Compteurs
\newcounter{seancenum}
\newcounter{seancenumD}
\setcounter{seancenum}{1}
\setcounter{seancenumD}{1}

% Configuration des en-têtes et pieds de page
\pagestyle{fancy}
\fancyhf{}
\ifthenelse{\boolean{prof}}{
    \fancyhead[L]{\textcolor{Couleur8}{\textbf{Corrections Automatismes - Classe de seconde - Version Professeur}}}
}{
    \fancyhead[L]{\textcolor{Couleur8}{\textbf{Corrections Devoirs - Séance 25 - Classe de seconde}}}
}
\fancyhead[R]{\textcolor{Couleur8}{\textbf{2025-2026}}}
\fancyfoot[L]{\textcolor{Couleur8}{Mme Bertail et Mme Pinto}}
\fancyfoot[R]{\textcolor{Couleur8}{\thepage}}
\renewcommand{\headrulewidth}{1pt}
\renewcommand{\headrule}{\hbox to\headwidth{\color{Couleur8}\leaders\hrule height \headrulewidth\hfill}}

% Environnements personnalisés
\newtcolorbox{seance}[1]{
    colback=Couleur6!10,
    colframe=Couleur1!65,
    fonttitle=\bfseries\large,
    title=Correction Séance \theseancenum #1,
    left=5mm,
    right=5mm,
    top=3mm,
    bottom=3mm,
    arc=0mm,
    after=\stepcounter{seancenum}
}

\newtcolorbox{seanceD}[1]{
    colback=Couleur6!5,
    colframe=Couleur6!45,
    fonttitle=\bfseries\large,
    title=Correction Devoirs - Séance \theseancenumD #1,
    left=5mm,
    right=5mm,
    top=3mm,
    bottom=3mm,
    arc=0mm,
    after=\stepcounter{seancenumD}
}

\begin{document}

\setcounter{seancenumD}{25}
\newpage
\begin{seanceD}{} %25

\begin{enumerate}
\item $\dfrac{5}{8}x+2x=\dfrac{5}{8}x+\dfrac{2\times8}{8}x=\dfrac{5}{8}x+\dfrac{16}{8}x=\dfrac{5+16}{8}x=$ ${\color[HTML]{f15929}\boldsymbol{\dfrac{21}{8}x}}$

\item On utilise l'égalité remarquable ${\color{red} a}^2-{\color{blue} b}^2=({\color{red} a}-{\color{blue} b})({\color{red} a}+{\color{blue} b})$ avec $a={\color{red} 4(3x+5)}$ et $b={\color{blue} 3(8x-1))}$.\\
  $\begin{aligned}16(3x+5)^2-9(8x-1)^2
  &=\underbrace{({\color{red} 4(3x+5)})^2-({\color{blue} 3(8x-1)})^2}_{a^2-b^2}\\
  &=\underbrace{[({\color{red} 4(3x+5)})-({\color{blue} 3(8x-1)})][({\color{red} 4(3x+5)})+({\color{blue} 3(8x-1)})]}_{(a-b)(a+b)}\\
  &= (12x+20-24x+3)(12x+20+24x-3)\\ &={\color[HTML]{f15929}\boldsymbol{(-12x+23)(36x+17)}}\end{aligned}$

\item Diminuer de $34\,\%$ revient à multiplier par $1 - 0{,}34$ (coefficient multiplicateur).\\
Si $V_I$ est le prix initial, on a : $ V_I \times \left(1 - 0{,}34\right)=160$.\\
Ainsi, le prix initial est donné par : ${\color[HTML]{f15929}\boldsymbol{\dfrac{160}{1 - 0{,}34}}}$.\\La bonne réponse est la réponse {\bfseries \color[HTML]{f15929}A}.

\item  $9$ $\text{ m}$ $= 0{,}9$ $\text{dam}$\\
L'aire du carré est : $0{,}9 \text{ dam}\times 0{,}9 \text{ dam} = {\color[HTML]{f15929}\boldsymbol{0{,}81\text{ dam}^2}}$.

\item On a $9=3^{2}$ donc :\\
    $
    \text{Le triple de  } 9^{50}  = 3 \times 9^{50} 
     = 3\times  \left(3^{2}\right)^{50} 
    =3\times  3^{100} 
    ={\color[HTML]{f15929}\boldsymbol{3^{101}}} 
$
   
\item $\dfrac{x-12}{x+1} > 0$\\$\bullet$ On commence par chercher les éventuelles valeurs interdites :\\$x +1 = 0$ \\$x +1 {\color[HTML]{f15929}\boldsymbol{-1}} = {\color[HTML]{f15929}\boldsymbol{-1}}$\\$x = -1$\\Le quotient est défini sur $\mathbb{R} \smallsetminus \{-1\}$.\\$\bullet$ On résout l'inéquation sur $\mathbb{R} \smallsetminus \{-1\}$.\\$x -12 > 0$ si et seulement si $x > 12$.\\$x+1 > 0$ si et seulement si $x > -1$.\\On peut donc en déduire le tableau de signes suivant :\\\begin{tikzpicture}[baseline, scale=0.75]
    \tkzTabInit[lgt=8,deltacl=0.8,espcl=3.5]{ $x$ / 2, $x-12$ / 2, $x+1$ / 2, $\dfrac{x-12}{x+1}$ / 4}{ $-\infty$, $-1$, $12$, $+\infty$}
	\tkzTabLine{  , -, t, -, z, +}
	\tkzTabLine{  , -, z, +, t, +}
	\tkzTabLine{  , +, d, -, z, +}
	\end{tikzpicture}
\\ L'ensemble de solutions de l'inéquation est $S = \left] -\infty ; -1 \right[ \cup \left] 12; +\infty \right[ $.
\end{enumerate}
\end{seanceD}

\end{document}